\subsection{对自由幺半群的应用}

引理2. 设 \(\mathbf{M}\) 是族 \((M_x)_{x\in \mathbf{X}}\) 的幺半和，该族由 \(\mathbf{M}_x = \mathbf{N}\) （对所有 \(x\in \mathbf{X}\) ）定义，并设 \(\phi_{x}\) 是 \(\mathbf{M}_x\) 到 \(\mathbf{M}\) 中的典范同态。映射 \(x\mapsto \phi_x(1)\) （从 \(\mathbf{X}\) 到 \(\mathbf{M}\) ）可扩张为一个同构 \(h\) （从 \(\operatorname {Mo}(\mathbf{X})\) 到 \(\mathbf{M}\) 上）。

设 \(h\) 是 \(\operatorname{Mo}(\mathbf{X})\) 到 \(\mathbf{M}\) 中的、以 \(h(x) = \phi_x(1)\) 为特征的同态。对每个整数 \(n \geqslant 0\) ，有 \(\phi_x(n) = \phi_x(1)^n = h(x)^n\) ，而且由于 \(\mathbf{M}\) 由 \(\bigcup_{x \in \mathbf{X}} \phi_x(\mathbf{N})\) 生成，它也由 \(h(\mathbf{X})\) 生成。因此 \(h\) 是满射。此外，对所有 \(x\) （属于 \(\mathbf{X}\) ），映射 \(n \mapsto x^n\) 是 \(\mathbf{N} = \mathbf{M}_x\) 到 \(\operatorname{Mo}(\mathbf{X})\) 中的同态；因此存在（第 3 号，命题 4）一个同态 \(h'\) ，它把 \(\mathbf{M}\) 映入 \(\operatorname{Mo}(\mathbf{X})\) ，使得 \(h'(\phi_x(n)) = x^n\) 对 \(x \in \mathbf{X}\) 和 \(n \in \mathbf{N}\) 成立；特别地， \(h'(h(x)) = x\) 对 \(x \in \mathbf{X}\) 成立，因此 \(h' \circ h\) 是 \(\operatorname{Mo}(\mathbf{X})\) 的恒等同态。于是 \(h\) 是单射。这样就证明了 \(h\) 是双射。

命题6. 设 \(w\) 是 \(\mathbf{Mo}(\mathbf{X})\) 的一个元素。

\begin{enumerate}
\def\labelenumi{(\alph{enumi})}
\tightlist
\item
  存在一个整数 \(n \geqslant 0\) 、元素 \(x_{\alpha}\) （属于 \(\mathbf{X}\) ）以及整数 \(m(\alpha) > 0\) （对于
\end{enumerate}

\(1 \leqslant \alpha \leqslant n\) ），使得 \(x_{\alpha} \neq x_{\alpha + 1}\) 对 \(1 \leqslant \alpha < n\) 成立，且 \(w = \prod_{\alpha = 1}^{n} x_{\alpha}^{m(\alpha)}\) 。序列 \((x_{\alpha}, m(\alpha))_{1 \leqslant \alpha \leqslant n}\) 由这些条件唯一确定。

\[
x _ {\beta} ^ {\prime}
\]

\[
m ^ {\prime} (\beta)
\]

\[
1 \leqslant \beta \leqslant p
\]

(b) \(w = \prod_{\beta=1}^{p} x'^{m'(\beta)}.\) 那么 \(p \geqslant n\) 。若 \(p = n\) ，则 \(x_{\beta}' = x_{\beta}\) 且 \(m'(\beta) = m(\beta)\) 对 \(1 \leqslant \beta \leqslant p\) 成立。

在引理2的记号下， \(\bar{h}^{1}(\phi_x(n)) = x^n\) 对 \(x \in \mathbf{X}\) 和 \(n \in \mathbf{N}\) 成立。命题 6 随后由第 3 号命题 5 得出。

